%---------------------------------------------------------------------
\chapter{Literature Survey and Existing Practice}
%---------------------------------------------------------------------

This chapter first describes how GI claims on handloom products are established and checked in
practice, and then reviews the technical literature on the agentic techniques used in LoomTrace. It
ends with a comparative summary and a statement of the gap that the project addresses.

\section{Existing Practice in the Sector}

\subsection{Legal basis of a GI claim}

Under the Geographical Indications of Goods (Registration and Protection) Act, 1999, a GI is
registered for a class of goods and a defined geographical area on the application of an association
of persons or producers. The application contains a \emph{statement of case} that describes the
product, the geographical area, the method of production, the materials and the link between the
product and its origin~\cite{giact1999}. After registration, producers within the area may apply to
be registered as \emph{authorised users}, and only the registered proprietor and authorised users have
the right to use the indication. The registry and its records are maintained by IP India~\cite{ipindia}.
The international basis for this protection is Section~3 of the TRIPS Agreement, which defines
geographical indications and obliges members to prevent misleading use of them~\cite{trips}.

\subsection{How a claim is checked today}

In practice, the verification of a GI claim at the point of sale is manual and fragmented. A buyer in
a shop or on an online marketplace sees a product description or label. To check the claim, the buyer
or a marketplace moderator must find the corresponding GI registration and its application number,
read the statement of case to learn the registered area and materials, and compare these with the
place and materials stated by the seller. Where producer details are provided, they must be checked
against cooperative or weaver records, which are not published centrally. The Ministry of Textiles
supports buyer confidence through separate labelling schemes, notably the Handloom Mark, which
certifies that a product is handwoven, and the India Handloom Brand, which certifies quality and
process standards~\cite{handloommark}. These labels are useful signals, but they do not certify that a
product falls within a particular GI, and they are not linked to the GI registry in a way that a buyer
can query.

Government publications provide some of the reference data needed for verification. The TRIFED GI
catalogue published through the Press Information Bureau (PIB) lists registered GI products marketed
by Tribes India together with short descriptions of their place of production and
materials~\cite{pibgi}. The state-wise results of the Fourth All India Handloom Census 2019--20,
published as a PIB annexure, give the number of handloom weavers and allied workers in each state and
union territory~\cite{pibcensus}. These sources are authoritative but are distributed as separate PDF
documents with different structures, and none of them is designed to answer the question ``is this
particular claim supported?''.

\subsection{Where current practice fails}

Three weaknesses of current practice motivate this project. First, the evidence is spread across
several documents that must be read together, so a complete check is slow and is rarely performed
for an individual purchase. Second, product names on labels and listings frequently differ from the
registered names, which makes simple string search unreliable. Third, when information is missing,
there is no defined route for escalation: the checker must either accept or reject the claim without
evidence, or abandon the check. A system that collects the evidence, explains which checks passed and
failed, and routes undecidable cases to a person addresses all three weaknesses.

\section{Related Technical Work}

The technical literature relevant to LoomTrace is organised below into four themes: reasoning and
acting with tools; retrieval-augmented generation; planning, reflection and multi-agent
collaboration; and the evaluation of systems built on language models.

\subsection{Reasoning and Acting with Tools}

Yao \emph{et al.}~\cite{yao2023react} proposed ReAct, in which a language model interleaves
free-text reasoning steps with actions such as search queries, so that it can gather information
before answering. Schick \emph{et al.}~\cite{schick2023toolformer} showed that a language model can
learn to decide when to call external tools such as calculators and search engines, and how to use
their results. Both works establish that grounding a model's output in tool results reduces reliance
on memorised knowledge. LoomTrace adopts the tool-use principle, since every fact used in a verdict is
returned by a typed tool with a source identifier. It does not, however, adopt free-form interleaving
of thought and action: the sequence of nodes is fixed by an explicit graph, so that every run can be
traced, bounded and checkpointed. This trade-off favours auditability over flexibility, which suits a
verification task whose steps are known in advance.

\subsection{Retrieval-Augmented Generation}

Lewis \emph{et al.}~\cite{lewis2020rag} introduced retrieval-augmented generation (RAG), which
conditions a generator on passages retrieved from an external corpus, and showed gains on
knowledge-intensive tasks. Gao \emph{et al.}~\cite{gao2023ragsurvey} survey later developments and
distinguish naive RAG, which retrieves once, from advanced and modular variants that rewrite queries,
re-rank results and retrieve iteratively. Two adaptive variants are particularly relevant. Self-RAG,
proposed by Asai \emph{et al.}~\cite{asai2024selfrag}, trains a model to decide when to retrieve and to
critique its own output with reflection tokens. Corrective RAG, proposed by Yan
\emph{et al.}~\cite{yan2024crag}, adds a lightweight evaluator that grades retrieved documents and
triggers corrective actions, such as a new search, when retrieval is weak. LoomTrace follows the
corrective pattern: its grader detects when the registry lookup or the document search did not
resolve the product, and the planner then re-queries with the canonical product name and a filter on
that product. For ranking, LoomTrace uses the BM25 function from the probabilistic relevance
framework of Robertson and Zaragoza~\cite{robertson2009bm25}, which is fast, transparent and effective
on short, keyword-rich official passages.

\subsection{Planning, Reflection and Multi-Agent Collaboration}

Wang \emph{et al.}~\cite{wang2023plansolve} proposed plan-and-solve prompting, in which a model first
devises a plan that divides a task into subtasks and then executes it, which reduces missed steps
compared with direct reasoning. Shinn \emph{et al.}~\cite{shinn2023reflexion} proposed Reflexion, in
which an agent evaluates its own attempt and uses verbal feedback to improve the next attempt. Wu
\emph{et al.}~\cite{wu2023autogen} presented AutoGen, a framework in which several specialised agents
converse to complete a task. A survey by Wang \emph{et al.}~\cite{wang2024survey} organises agents built
on language models by their profile, memory, planning and action modules. LoomTrace draws on each of
these ideas in a constrained form. A planner decides which checks a claim requires; a grader evaluates
the evidence before any verdict is produced and can send control back to the planner; and the
gathering of evidence and the judging of evidence are assigned to two separate nodes, a retriever agent
and a verifier agent, with the graph acting as the supervisor that routes between them.

\subsection{Evaluation of LLM-Based Systems}

Zheng \emph{et al.}~\cite{zheng2023judge} studied the use of strong language models as judges of other
models' outputs and reported high agreement with human preferences, while documenting position,
verbosity and self-enhancement biases. Ji \emph{et al.}~\cite{ji2023hallucination} survey hallucination
in natural language generation and emphasise that fluency is not evidence of faithfulness. Liu
\emph{et al.}~\cite{liu2024lost} show that models use information at the beginning and end of long
contexts better than information in the middle. For LoomTrace, these findings have two consequences.
First, because each verdict is one of three discrete labels and each check has a deterministic expected
outcome, the system is evaluated with exact-match metrics against hand-assigned labels rather than with
a model judge. Second, the design keeps every model call narrow and short and assigns the factual checks
to deterministic rules, which is consistent with the known weaknesses of small models on long inputs.

\section{Comparative Summary}

Table~\ref{tab:lit} compares the most relevant approaches with LoomTrace on four properties: whether
the approach calls external tools, whether retrieval is adaptive (graded and repeated when weak),
whether it includes a human escalation path, and whether every decision is accompanied by citations to
sources.

\begin{table}[H]
\centering
\caption{Comparison of related approaches with the proposed system}
\label{tab:lit}
\small
\begin{tabularx}{\linewidth}{@{}>{\raggedright\arraybackslash}p{3.2cm} >{\raggedright\arraybackslash}X c c c c@{}}
\toprule
\textbf{Work} & \textbf{Key idea} & \textbf{Tools} & \textbf{Adapt. RAG} & \textbf{HITL} & \textbf{Cited} \\ \midrule
ReAct~\cite{yao2023react} & Interleaved reasoning and acting & \checkmark & -- & -- & -- \\
Toolformer~\cite{schick2023toolformer} & Self-taught tool calling & \checkmark & -- & -- & -- \\
RAG~\cite{lewis2020rag} & Retrieval-conditioned generation & -- & -- & -- & \checkmark \\
Self-RAG~\cite{asai2024selfrag} & Learned retrieval and self-critique & -- & \checkmark & -- & \checkmark \\
CRAG~\cite{yan2024crag} & Retrieval grading and correction & -- & \checkmark & -- & \checkmark \\
Reflexion~\cite{shinn2023reflexion} & Verbal self-feedback across attempts & \checkmark & -- & -- & -- \\
Manual GI check~\cite{giact1999,ipindia} & Person reads registry and records & -- & -- & \checkmark & \checkmark \\
\textbf{LoomTrace (ours)} & Graph agent with graded re-query and review queue & \checkmark & \checkmark & \checkmark & \checkmark \\
\bottomrule
\end{tabularx}
\end{table}

\section{Research Gap}

As Table~\ref{tab:lit} shows, the general agent techniques reviewed provide tool use, adaptive
retrieval or self-critique individually, but they are not applied to GI verification, and none of them
defines an escalation path to a human reviewer. Manual verification provides human judgement and
citations but does not scale and has no structured way of handling missing information. No reviewed
approach combines retrieval over official GI text, structured checks against registry, census and
producer records, a bounded re-planning loop, and a human review workflow in a single traceable
system for handloom GI claims; LoomTrace is designed to address this gap.
